% =============================================================================
%  InsaCloud — Script de soutenance
%  Compilation :  tectonic -X compile soutenance.tex
% =============================================================================
\documentclass[a4paper,11pt,oneside]{report}
% =============================================================================
%  Préambule du guide technique InsaCloud
% =============================================================================
\usepackage{fontspec}
% Noms de fichiers plutôt que noms système : la police vient du paquet
% tectonic, le document se recompile donc à l'identique sur n'importe quel poste.
\setmainfont{lmroman10-regular.otf}[
  BoldFont=lmroman10-bold.otf, ItalicFont=lmroman10-italic.otf,
  BoldItalicFont=lmroman10-bolditalic.otf]
\setsansfont{lmsans10-regular.otf}[
  BoldFont=lmsans10-bold.otf, ItalicFont=lmsans10-oblique.otf]
\setmonofont{lmmono10-regular.otf}[
  BoldFont=lmmonolt10-bold.otf, ItalicFont=lmmono10-italic.otf, Scale=0.92]
\usepackage[french]{babel}
\usepackage{graphicx,xcolor,booktabs,array,tabularx,longtable,multirow}
\usepackage{tikz}
\usetikzlibrary{positioning,arrows.meta,shapes.geometric,calc,fit,backgrounds,
                decorations.pathreplacing,patterns,shadows.blur}
\usepackage{fancyhdr,enumitem,microtype}
\usepackage[colorlinks=true,linkcolor=insared,urlcolor=insared,citecolor=insared]{hyperref}
\usepackage{caption}
\usepackage{listings}
\usepackage{tcolorbox}
\tcbuselibrary{skins,breakable}
\usepackage{titlesec}
\usepackage{emptypage}

% --- Couleurs ----------------------------------------------------------------
\definecolor{insared}{RGB}{200,16,46}
\definecolor{grisfond}{RGB}{247,248,250}
\definecolor{grisligne}{RGB}{216,220,226}
\definecolor{gristexte}{RGB}{91,97,107}
\definecolor{vertok}{RGB}{26,127,75}
\definecolor{vertfond}{RGB}{238,247,241}
\definecolor{orangepartiel}{RGB}{178,106,5}
\definecolor{orangefond}{RGB}{254,246,231}
\definecolor{rougeko}{RGB}{198,40,40}
\definecolor{bleufond}{RGB}{238,243,250}
\definecolor{bleutrait}{RGB}{43,91,158}

% --- Mise en page ------------------------------------------------------------
\captionsetup{font=small,labelfont={bf,color=insared}}
\setlist[itemize]{leftmargin=1.3em,itemsep=2pt,topsep=3pt,parsep=0pt}
\setlist[enumerate]{leftmargin=1.5em,itemsep=2pt,topsep=3pt,parsep=0pt}
\setlist[description]{leftmargin=1.3em,itemsep=3pt,topsep=3pt,parsep=0pt}

% LaTeX refuse par defaut qu'une figure occupe plus de 70 % d'une page et
% exige 20 % de texte : avec beaucoup de schemas, cela repousse les flottants
% et laisse des pages a moitie vides. On relache ces seuils.
\renewcommand{\topfraction}{0.92}
\renewcommand{\bottomfraction}{0.85}
\renewcommand{\textfraction}{0.06}
\renewcommand{\floatpagefraction}{0.75}
\setcounter{topnumber}{3}
\setcounter{bottomnumber}{2}
\setcounter{totalnumber}{4}

\pagestyle{fancy}
\fancyhf{}
\fancyhead[L]{\small\textcolor{gristexte}{\leftmark}}
\fancyhead[R]{\small\textcolor{gristexte}{InsaCloud}}
\fancyfoot[C]{\small\thepage}
\renewcommand{\headrulewidth}{0.4pt}
\renewcommand{\headrule}{\hbox to\headwidth{\color{grisligne}\leaders\hrule height \headrulewidth\hfill}}

\titleformat{\chapter}[display]
  {\normalfont\huge\bfseries}
  {\textcolor{insared}{\normalsize\MakeUppercase{\chaptertitlename} \thechapter}}
  {6pt}{\Huge}
\titlespacing*{\chapter}{0pt}{-20pt}{24pt}

% --- Extraits de code --------------------------------------------------------
\lstdefinestyle{insacloud}{
  basicstyle=\ttfamily\footnotesize,
  backgroundcolor=\color{grisfond},
  frame=single, rulecolor=\color{grisligne},
  breaklines=true, breakatwhitespace=false,
  postbreak=\mbox{\textcolor{gristexte}{$\hookrightarrow$}\space},
  showstringspaces=false,
  keywordstyle=\color{insared}\bfseries,
  commentstyle=\color{gristexte}\itshape,
  stringstyle=\color{vertok},
  numberstyle=\tiny\color{gristexte},
  xleftmargin=6pt, framexleftmargin=6pt, aboveskip=8pt, belowskip=8pt,
  extendedchars=true, inputencoding=utf8,
  literate=%
    {é}{{\'e}}1 {è}{{\`e}}1 {ê}{{\^e}}1 {ë}{{\"e}}1
    {à}{{\`a}}1 {â}{{\^a}}1 {ä}{{\"a}}1
    {î}{{\^i}}1 {ï}{{\"i}}1
    {ô}{{\^o}}1 {ö}{{\"o}}1
    {û}{{\^u}}1 {ù}{{\`u}}1 {ü}{{\"u}}1
    {ç}{{\c c}}1 {É}{{\'E}}1 {È}{{\`E}}1 {À}{{\`A}}1 {Ç}{{\c C}}1
    {œ}{{\oe}}1 {Œ}{{\OE}}1 {æ}{{\ae}}1
    {—}{{---}}1 {–}{{--}}1 {’}{{'}}1 {‘}{{'}}1
    {«}{{\guillemotleft}}1 {»}{{\guillemotright}}1
    {·}{{\textperiodcentered}}1 {→}{{$\rightarrow$}}1 {←}{{$\leftarrow$}}1
    {✔}{{\checkmark}}1 {°}{{\textdegree}}1 {€}{{\texteuro}}1
}
\lstset{style=insacloud}

\lstdefinelanguage{yamlplus}{
  keywords={true,false,null,name,hosts,become,roles,tasks,vars,when,loop,with_items,
            state,mode,owner,group,path,src,dest,notify,handlers,services,image,
            environment,volumes,networks,ports,depends_on,build,command,restart},
  keywordstyle=\color{insared}\bfseries,
  sensitive=false,
  comment=[l]{\#},
  morestring=[b]",
  morestring=[b]',
}

\lstdefinelanguage{dockerfile}{
  keywords={FROM,RUN,CMD,LABEL,EXPOSE,ENV,ADD,COPY,ENTRYPOINT,VOLUME,USER,
            WORKDIR,ARG,ONBUILD,STOPSIGNAL,HEALTHCHECK,SHELL,AS},
  keywordstyle=\color{insared}\bfseries,
  sensitive=true,
  comment=[l]{\#},
  morestring=[b]",
}

% --- Encadrés pédagogiques ---------------------------------------------------
\newtcolorbox{notion}[1]{
  breakable, enhanced, colback=bleufond, colframe=bleutrait,
  boxrule=0pt, leftrule=3pt, arc=2pt, left=8pt, right=8pt, top=6pt, bottom=6pt,
  fonttitle=\bfseries\small, title={Notion — #1},
  coltitle=bleutrait, titlerule=0pt, colbacktitle=bleufond,
  attach title to upper=\par\smallskip
}

\newtcolorbox{piege}[1]{
  breakable, enhanced, colback=orangefond, colframe=orangepartiel,
  boxrule=0pt, leftrule=3pt, arc=2pt, left=8pt, right=8pt, top=6pt, bottom=6pt,
  fonttitle=\bfseries\small, title={Piège — #1},
  coltitle=orangepartiel, titlerule=0pt, colbacktitle=orangefond,
  attach title to upper=\par\smallskip
}

\newtcolorbox{retenir}{
  breakable, enhanced, colback=vertfond, colframe=vertok,
  boxrule=0pt, leftrule=3pt, arc=2pt, left=8pt, right=8pt, top=6pt, bottom=6pt,
  fonttitle=\bfseries\small, title={À retenir},
  coltitle=vertok, titlerule=0pt, colbacktitle=vertfond,
  attach title to upper=\par\smallskip
}

% --- Raccourcis --------------------------------------------------------------
\newcommand{\ok}{\textcolor{vertok}{\textbf{Conforme}}}
\newcommand{\partiel}{\textcolor{orangepartiel}{\textbf{Partiel}}}
\newcommand{\plus}{\textcolor{insared}{\textbf{Au-delà}}}
\newcommand{\fichier}[1]{\textcolor{insared}{\texttt{#1}}}
\newcommand{\cmd}[1]{\texttt{#1}}

% Les paragraphes techniques enchaînent de longues suites en \texttt{} que TeX
% ne peut pas couper : sans marge de manœuvre, elles débordent dans la marge.
\emergencystretch=3em
\hbadness=6000
\tolerance=2000


\usepackage{tcolorbox}
\tcbuselibrary{skins,breakable}

% --- Encadrés propres à ce document ------------------------------------------
\newtcolorbox{adire}{
  breakable, enhanced, colback=white, colframe=insared,
  boxrule=0pt, leftrule=3.5pt, arc=2pt, left=10pt, right=10pt, top=7pt, bottom=7pt,
  fontupper=\itshape
}

\newtcolorbox{afaire}[1]{
  breakable, enhanced, colback=bleufond, colframe=bleutrait,
  boxrule=0pt, leftrule=3.5pt, arc=2pt, left=10pt, right=10pt, top=6pt, bottom=6pt,
  fonttitle=\bfseries\small, title={À faire — #1},
  coltitle=bleutrait, titlerule=0pt, colbacktitle=bleufond,
  attach title to upper=\par\smallskip
}

\newcommand{\minute}[1]{%
  \vspace{4pt}\noindent{\color{insared}\rule{3pt}{14pt}}\hspace{6pt}%
  {\large\bfseries #1}\par\vspace{3pt}}

\newcommand{\fonc}[1]{%
  \par\vspace{3pt}\noindent{\small\color{insared}$\rightarrow$\ \textbf{#1}}\par}

\newcommand{\capture}[2]{%
  \begin{figure}[htbp]\centering
  \includegraphics[width=#2\textwidth,frame]{captures/#1}
  \end{figure}}

\usepackage[export]{adjustbox}

\begin{document}

% --- Couverture --------------------------------------------------------------
\begin{titlepage}
\centering
\vspace*{2.5cm}

{\Huge\bfseries InsaCloud\par}
\vspace{6pt}
{\color{insared}\rule{4cm}{2pt}\par}
\vspace{18pt}
{\LARGE Script de soutenance\par}
\vspace{8pt}
{\large 15 minutes · trois intervenants · démonstration en direct\par}

\vspace{2cm}
\begin{minipage}{0.84\textwidth}
\color{gristexte}\small
À chaque geste de la démonstration, celui qui parle nomme \textbf{l'outil
employé} et \textbf{la fonctionnalité du cahier des charges} qu'il réalise. La
prof doit pouvoir cocher sa grille sans avoir à nous le demander.

\medskip
Toutes les commandes de ce document ont été exécutées sur la plateforme. Toutes
les captures d'écran viennent de l'application telle qu'elle tourne, et les
sorties reproduites sont les vraies.
\end{minipage}

\vfill
{\color{gristexte}\small
Projet 3 --- \textit{Outils de déploiement de plateformes}\\
INSA Centre-Val de Loire, STI 4A FISE, 2026/2027 --- Dr.~Nadia Fettah\\[8pt]
\today\par}
\end{titlepage}

\pagenumbering{roman}
\setcounter{tocdepth}{0}
\tableofcontents
\cleardoublepage
\pagenumbering{arabic}

% =============================================================================
\chapter*{Avant d'entrer dans la salle}
\addcontentsline{toc}{chapter}{Avant d'entrer dans la salle}
\markboth{PRÉPARATION}{}

À faire \textbf{vingt minutes} avant, pas deux.

\begin{lstlisting}[language=bash]
colima start --cpu 4 --memory 6 --disk 40
cd ~/Developer/insacloud && docker compose -p insacloud up -d
\end{lstlisting}

\paragraph{Vérifier les trois conteneurs.} Les trois doivent être
\cmd{healthy}. S'il en manque un, ne pas entrer : relancer.

\begin{lstlisting}[language=bash]
docker compose -p insacloud ps
\end{lstlisting}

\lstinputlisting[basicstyle=\ttfamily\scriptsize]{sorties/pile.txt}

\paragraph{Vérifier les six images.} Sans elles, le mode bureau n'apparaîtra
pas dans le formulaire --- et ce n'est pas un bug, c'est voulu : le catalogue
ne propose pas ce qu'il ne peut pas livrer.

\begin{lstlisting}[language=bash]
docker images --filter 'reference=insacloud_*'
\end{lstlisting}

\lstinputlisting[basicstyle=\ttfamily\scriptsize]{sorties/images.txt}

\paragraph{Préparer les onglets}, dans cet ordre :

\begin{enumerate}
\item \cmd{http://127.0.0.1:8088} --- connecté, tableau de bord \textbf{vide}
\item \cmd{https://github.com/aziz3r/insacloud/actions} --- la chaîne au vert
\item un terminal plein écran dans \cmd{\textasciitilde/Developer/insacloud}
\item le \cmd{README.md} du dépôt, prêt à montrer les schémas
\end{enumerate}

Un \textbf{second terminal}, pour la panne du nœud, avec les deux commandes
déjà tapées mais non validées : \cmd{docker logs -f insacloud-faucheur} dans
l'un, \cmd{docker stop insacloud-noeud1} dans l'autre.

Et un cinquième onglet sur \cmd{http://127.0.0.1:8088/workers} : l'état du
parc, que les agents alimentent eux-mêmes.

\begin{retenir}
Personne ne tape une commande qu'il n'a pas répétée. Chaque commande de ce
document a été exécutée au moins une fois avant la soutenance.
\end{retenir}

\section*{Répartition}

\begin{center}
\small
\begin{tabularx}{\textwidth}{@{}llX@{}}
\toprule
\textbf{Qui} & \textbf{Quand} & \textbf{Rôle} \\
\midrule
Personne 1 & 0:00 -- 3:30 & Le problème, l'architecture, l'infrastructure \\
Personne 2 & 3:30 -- 9:00 & Le parcours utilisateur, en direct \\
Personne 3 & 9:00 -- 13:30 & Les scénarios de l'énoncé, la sécurité, la chaîne \\
Personne 1 & 13:30 -- 15:00 & Conformité, ajouts, conclusion \\
\bottomrule
\end{tabularx}
\end{center}

% =============================================================================
\chapter{Personne 1 — Le problème et l'infrastructure}
\markboth{PERSONNE 1 — 0:00 À 3:30}{}

\minute{0:00 — Ouvrir (30 s)}

\begin{adire}
InsaCloud est un mini-fournisseur de cloud. Un utilisateur s'inscrit, choisit
une distribution Linux et une durée, et obtient en quinze secondes une machine
à lui, avec un accès root. À l'échéance, elle disparaît toute seule.

\medskip
Derrière ce clic il y a quatre machines, un ordonnanceur, un démon de
nettoyage, un coffre à secrets et onze contrôles automatiques. C'est ce que
nous allons dérouler.
\end{adire}

\minute{0:30 — Le choix fondateur (45 s)}

\begin{afaire}{montrer}
Le schéma d'architecture, dans le \cmd{README.md} du dépôt.
\end{afaire}

\begin{adire}
Première décision : des \textbf{conteneurs}, pas des machines virtuelles. Une
VM démarre en trente secondes et réserve des centaines de mégaoctets rien que
pour son noyau. Un conteneur réutilise le noyau de l'hôte : il démarre en moins
d'une seconde. Pour une location \textbf{à la minute}, c'est décisif --- et
l'énoncé le tranche lui-même : « des conteneurs dans notre cas ».

\medskip
Deuxième décision : séparer \textbf{qui décide} de \textbf{qui exécute}. Le
contrôleur porte le site, la base et l'ordonnanceur. Il n'exécute aucune machine
louée et n'a même pas de démon Docker, seulement le client. Les trois workers
portent Docker et les conteneurs. Un utilisateur qui s'échapperait de son
conteneur tomberait sur un worker : pas de base, pas de clés, pas de code.
\end{adire}

\minute{1:15 — Comment le contrôleur pilote les workers (45 s)}

\begin{afaire}{taper}
\begin{lstlisting}[language=bash]
docker -H ssh://insacloud@worker1 ps
\end{lstlisting}
En mode compose, montrer plutôt \fichier{2\_webapp/workers.py}, fonction
\cmd{docker\_command}.
\end{afaire}

\begin{adire}
On aurait pu exposer l'API Docker sur le réseau, port 2375. C'est simple et
c'est une très mauvaise idée : qui parle au démon Docker d'une machine en est
root. Une erreur de pare-feu et toute la plateforme tombe.

\medskip
On passe donc par \textbf{SSH}, que le client Docker sait faire nativement.
Aucun port Docker n'est ouvert, l'authentification est celle de SSH par clé, et
le compte de service n'a pas de mot de passe.
\end{adire}

\fonc{Fonctionnalité 4 (conteneurisation) · Fonctionnalité 14 (sécurité)}

\minute{2:00 — L'infrastructure comme code (1 min)}

\begin{afaire}{montrer}
Le \fichier{Vagrantfile}, puis \fichier{3\_ansible/site.yml}.
\end{afaire}

\begin{adire}
Rien n'est installé à la main. \textbf{Vagrant} décrit les quatre machines
virtuelles --- un contrôleur en .10, trois workers en .11, .12, .13 --- avec
leur mémoire, leurs cœurs, leur réseau privé. Le fichier gère quatre
\emph{providers}, dont VMware, parce que VirtualBox ne tourne pas sur une puce
Apple.

\medskip
\textbf{Ansible} les configure : neuf rôles, quatre playbooks. Pare-feu, clé de
pilotage, Docker, certificats TLS, images, agent, proxy inverse, application.

\medskip
Et surtout : le playbook est \textbf{idempotent}. Le relancer sur un serveur
déjà configuré ne change rien. La preuve se lit dans le compte rendu : au second
passage, \cmd{changed=0} sur les quatre nœuds. Ce n'est pas « ça ne fait rien la
deuxième fois » : c'est « il vérifie l'état et n'agit que s'il diffère ».
\end{adire}

\begin{adire}
Et il faut le dire ici, sans attendre la question : \textbf{aucun hyperviseur
ne démarre sur ce poste}. VMware n'a pas ses interfaces réseau --- son
extension système n'est pas installée --- et le support ARM de VirtualBox fait
tuer la machine virtuelle au démarrage. Le \fichier{Vagrantfile} et les
playbooks sont donc montrés et expliqués, pas exécutés devant vous.

\medskip
En revanche, la partie que l'énoncé teste vraiment --- un nœud qui tombe, ses
machines qui repartent ailleurs --- tourne \textbf{en direct} dans un instant,
sur deux nœuds en conteneurs.
\end{adire}

\begin{afaire}{montrer la commande, sans la lancer}
\begin{lstlisting}[language=bash]
cd 3_ansible
ansible-playbook -i inventaire_vagrant.py site.yml --vault-password-file .vault_pass
\end{lstlisting}
\end{afaire}

\fonc{Fonctionnalité 5 (Vagrant) · Fonctionnalité 6 (Ansible)}

\minute{3:00 — Les secrets (30 s)}

\begin{afaire}{montrer}
\fichier{3\_ansible/group\_vars/insacloud/vault.yml} --- illisible à l'écran,
c'est le but.
\end{afaire}

\begin{adire}
Les secrets vivent dans un \textbf{Ansible Vault}, chiffré en AES-256. Ce
fichier est versionné dans Git tel quel : sans la phrase de passe, il ne dit
rien. Il contient la clé de session Flask, la clé du coffre, et le jeton des
agents --- trois clés distinctes, jamais la même.
\end{adire}

\fonc{Fonctionnalité 6 · Fonctionnalité 14}

\begin{center}\small\itshape
Transition : « \textit{[Prénom 2]} va vous montrer ce que ça donne du côté de
l'utilisateur. »
\end{center}

% =============================================================================
\chapter{Personne 2 — Le parcours utilisateur}
\markboth{PERSONNE 2 — 3:30 À 9:00}{}

\begin{piege}{La règle de cette partie}
Tout se fait \textbf{en cliquant}, pas en racontant. On parle \emph{pendant} que
ça charge --- c'est là que se gagne la note.
\end{piege}

\minute{3:30 — Inscription et connexion (40 s)}

\capture{01-connexion.jpg}{0.62}

\begin{afaire}{faire}
Créer un compte en direct, puis se connecter.
\end{afaire}

\begin{adire}
L'application est en \textbf{Flask}, servie par \textbf{Gunicorn} avec quatre
processus ouvriers, derrière \textbf{nginx} qui termine le TLS. Flask seul
traite une requête à la fois et le dit lui-même : à ne pas utiliser en
production.

\medskip
La politique de mot de passe est vérifiée \textbf{côté serveur}. Un contrôle
fait seulement dans le navigateur se contourne en coupant JavaScript. Les mots
de passe de compte sont \textbf{hachés} en pbkdf2-sha256 : opération à sens
unique, on ne peut pas les retrouver. Et cinq échecs de connexion verrouillent
le compte.
\end{adire}

\fonc{Fonctionnalité 2 (API Web Flask) · Fonctionnalité 14}

\minute{4:10 — Le catalogue (30 s)}

\capture{02-tableau-de-bord.jpg}{0.92}

\begin{adire}
Trois distributions --- Ubuntu, Debian, Alpine --- en deux modes, soit
\textbf{six images} construites depuis \textbf{quatre Dockerfiles}. Le mode
bureau s'active par un simple \cmd{--build-arg DESKTOP=1}.

\medskip
Le catalogue n'est pas écrit en dur : il est \textbf{en base}, et l'application
interroge Docker avant d'afficher un choix. Si une image n'est pas construite,
la distribution n'apparaît pas --- on ne propose pas ce qu'on ne peut pas
livrer.
\end{adire}

\fonc{Fonctionnalité 3 (base de données) · Fonctionnalité 4}

\minute{4:40 — Louer (50 s)}

\begin{afaire}{faire}
Ubuntu · Bureau graphique · 10 minutes · \textbf{Louer la machine}.
Parler pendant les quinze secondes de démarrage.
\end{afaire}

\begin{adire}
Ce qui part côté serveur, c'est exactement la commande imposée par l'énoncé :

\medskip
\cmd{docker run -d --restart=always -p <port>:22 <image>}

\medskip
On la pilote par \textbf{\cmd{subprocess}}, pas par un SDK --- comme ça chaque
option se lit telle quelle dans le code. Et jamais avec \cmd{shell=True} : les
arguments sont passés en \textbf{liste}, donc aucun interpréteur ne relit la
chaîne. Sinon un utilisateur nommerait sa machine \cmd{x; rm -rf /}.

\medskip
S'y ajoutent le durcissement : \cmd{--cap-drop ALL} puis seulement les dix
capacités dont \cmd{sshd} a besoin, \cmd{no-new-privileges}, une limite de
processus, de CPU, et la mémoire --- \textbf{réservation et limite}, comme le
demande le cours.

\medskip
Le port vient de la plage 8000--9000, et son unicité est garantie par un
\textbf{index unique partiel} en base, limité aux machines actives. Deux
requêtes simultanées ne peuvent pas prendre le même port.
\end{adire}

\fonc{Fonctionnalité 9 (service de location) · Fonctionnalité 10 (déploiement
des instances) · Fonctionnalité 8 (ordonnanceur)}

\minute{5:30 — Le secret montré une fois (40 s)}

\capture{03-mot-de-passe-unique.jpg}{0.92}

\begin{adire}
Le mot de passe root est tiré au sort, seize caractères, environ 90 bits. Il est
affiché \textbf{une seule fois}, puis chiffré avec \textbf{Fernet} avant d'être
écrit en base --- clé distincte de celle des sessions.

\medskip
Une erreur qu'on a commise au début : on l'avait mis dans le cookie de session.
Or une session Flask est \textbf{signée, pas chiffrée} : n'importe qui pouvait
la décoder en base64. On ne stocke plus que l'identifiant de la machine, et le
secret est relu en base.
\end{adire}

\begin{afaire}{faire}
Recharger la page : le mot de passe a disparu. Puis ressaisir le mot de passe
du compte --- le coffre s'ouvre cinq minutes.
\end{afaire}

\capture{06-coffre-ouvert.jpg}{0.92}

\fonc{Fonctionnalité 14 (protection des secrets)}

\minute{6:10 — Les trois accès (1 min 20)}

\begin{afaire}{faire, dans cet ordre}
\textbf{1. SSH} --- copier la commande affichée, se connecter depuis le
terminal.\\
\textbf{2. Terminal web} --- ouvrir le lien.\\
\textbf{3. Bureau graphique} --- ouvrir, saisir le mot de passe VNC.
\end{afaire}

\begin{adire}
\textbf{SSH}, c'est l'accès demandé par l'énoncé. Par mot de passe, ou par clé
si l'utilisateur a enregistré sa clé publique --- dans ce cas la machine
n'accepte \textbf{plus que} la clé.
\end{adire}

\capture{05-terminal-web.jpg}{0.78}

\begin{adire}
Le \textbf{terminal web}, c'est \cmd{ttyd}. Attention : le port ne donne pas un
shell ouvert, il lance \cmd{login}. Il faut s'authentifier.
\end{adire}

\capture{04-bureau-xfce.jpg}{0.78}

\begin{adire}
Le \textbf{bureau}, c'est \cmd{Xvnc}, \cmd{XFCE} et \cmd{noVNC}. Cinq processus
cohabitent dans ce conteneur : sshd, ttyd, le serveur d'affichage, le bureau et
la passerelle web. Or un conteneur n'exécute qu'un processus --- d'où
\textbf{supervisord} en PID 1, qui les démarre et relance celui qui meurt.
\end{adire}

\fonc{Fonctionnalité 11 (accès SSH) · Fonctionnalité 4}

\minute{7:30 — L'expiration (1 min)}

\begin{afaire}{faire}
Louer une machine \textbf{1 minute}. Pendant qu'elle vit :
\begin{lstlisting}[language=bash]
docker logs -f insacloud-faucheur
\end{lstlisting}
\end{afaire}

\begin{adire}
Personne ne viendra cliquer sur « supprimer » à trois heures du matin. Un
service séparé, qu'on appelle le \textbf{Faucheur}, lit la base toutes les dix
secondes et détruit ce qui a expiré.

\medskip
Il fait aussi de la \textbf{réconciliation} : la base dit ce qui \emph{devrait}
exister, Docker ce qui existe \emph{vraiment}. Un conteneur orphelin est
supprimé, une instance sans conteneur passe à « arrêtée ». Avec une règle
essentielle : un nœud injoignable est \textbf{ignoré}, jamais interprété comme
vide --- sinon on marquerait arrêtées des machines qui tournent très bien.

\medskip
Et il reconnaît ses machines à une \textbf{étiquette Docker}, jamais à leur nom.
Au début il les reconnaissait par préfixe --- et il a fini par détruire les
conteneurs de la plateforme elle-même.
\end{adire}

\begin{afaire}{montrer}
La machine disparaît du tableau de bord, le quota repasse à zéro.
Puis les étiquettes qui lui donnent son identité :
\end{afaire}

\lstinputlisting[basicstyle=\ttfamily\scriptsize]{sorties/labels.txt}

\fonc{Fonctionnalité 12 (gestion de l'expiration)}

\minute{8:30 — L'ajout : la sauvegarde du travail (30 s)}

\capture{07-archives.jpg}{0.92}

\begin{adire}
Ça, l'énoncé ne le demande pas. Détruire un conteneur détruit sa couche
inscriptible : l'utilisateur perd tout. Juste avant le \cmd{docker rm -f}, on
fait un \cmd{tar} de \cmd{/root} et \cmd{/home} qu'on lit au fil de l'eau, et
l'utilisateur le télécharge pendant sept jours.

\medskip
Avec une règle qui prime sur tout : l'archivage \textbf{ne bloque jamais} la
destruction. Sinon une machine expirée sur un nœud en panne resterait en vie
indéfiniment.
\end{adire}

\fonc{Au-delà de l'énoncé}

\begin{center}\small\itshape
Transition : « \textit{[Prénom 3]} va vous montrer ce qui se passe quand ça
casse. »
\end{center}

% =============================================================================
\chapter{Personne 3 — Scénarios, sécurité, chaîne}
\markboth{PERSONNE 3 — 9:00 À 13:30}{}

\minute{9:00 — Le parc qui se gère seul (50 s)}

\begin{afaire}{taper}
\begin{lstlisting}[language=bash]
curl -s -H 'Accept: application/json' http://127.0.0.1:8088/workers | jq
\end{lstlisting}
\end{afaire}

\begin{adire}
Le contrôleur ne contient \textbf{aucune adresse de worker}. Chaque nœud porte
un petit \textbf{agent}, écrit en bibliothèque standard seulement --- pas de
dépendance à installer sur les workers. Au démarrage il se présente, puis il
signale sa présence toutes les trente secondes.

\medskip
Trois états : \cmd{AVAILABLE}, \cmd{BUSY} quand le nœud est plein,
\cmd{OFFLINE} après quatre-vingt-dix secondes de silence. La distinction
compte : un nœud saturé est en bonne santé, il ne faut surtout pas migrer ses
machines.

\medskip
Conséquence concrète : pour ajouter un quatrième worker, on l'allume. C'est
tout.
\end{adire}

\fonc{Fonctionnalité 7 (Worker Agent) · Fonctionnalité 8 (ordonnanceur)}

\minute{9:50 — Le Test 6 de l'énoncé : la panne (1 min 40)}

\begin{retenir}
Cette manipulation se fait \textbf{en direct}, et elle a été répétée : trente
et une secondes entre la coupure du nœud et la machine de nouveau en service.
\end{retenir}

\begin{afaire}{faire --- à gauche suivre le Faucheur, à droite couper le nœud}
\begin{lstlisting}[language=bash]
docker logs -f insacloud-faucheur      # terminal de gauche
docker stop insacloud-noeud1           # terminal de droite
\end{lstlisting}
\end{afaire}

\begin{adire}
L'énoncé décrit précisément ce test : on éteint un nœud, et on vérifie qu'il
passe \cmd{OFFLINE} et que l'instance repart sur un autre.

\medskip
L'agent cesse de battre. Premier message du journal : le nœud est injoignable
pour la réconciliation --- \textbf{et rien n'est détruit}. C'est volontaire :
un nœud muet n'est pas interprété. On ne sait pas s'il est mort ou seulement
coupé du réseau, et détruire sur un doute ferait perdre des machines qui
tournent très bien.

\medskip
Puis le délai de détection expire --- trente secondes ici, quatre-vingt-dix en
production --- et le contrôleur déclare le nœud hors ligne. Il recrée alors ses
machines sur un nœud disponible.

\medskip
Ce qui compte, c'est ce que le code \textbf{ne} touche \textbf{pas} : la table
des locations. C'est la décision qui structure toute la base --- on a séparé
\cmd{Instance}, qui décrit un conteneur, de \cmd{Rental}, qui décrit
l'engagement envers l'utilisateur. Si on avait fondu les deux en une table, plus
court à écrire, la reprise serait \textbf{impossible} : on ne pourrait pas
recréer la machine ailleurs sans toucher à la durée réservée.

\medskip
L'utilisateur garde donc sa location et sa durée ; il reçoit un nouveau port et
un nouveau mot de passe. Et un compteur \cmd{migrations} passe à 1 : c'est la
preuve observable, celle que vérifie notre test d'intégration.

\medskip
Une limite qu'on assume : la reprise ne rend \textbf{pas les fichiers}. Ils
vivaient dans la couche inscriptible, sur un nœud devenu injoignable. Sans
stockage partagé, seul le service repart. Un volume répliqué changerait la
nature du projet.
\end{adire}

\begin{piege}{Ce qui est simulé, et ce qui ne l'est pas}
Si la question vient --- et elle peut venir --- la réponse est courte et
vérifiable. Les deux nœuds sont des conteneurs, pas des machines virtuelles :
ils partagent le démon Docker de l'hôte au lieu d'en avoir un chacun. Tout le
reste est réel : le contrôleur les pilote par \cmd{ssh://insacloud@noeud1},
leur agent s'inscrit tout seul, le battement de cœur passe par l'API, la
bascule \cmd{OFFLINE} est celle du code, la recréation ailleurs aussi.

\medskip
Et cette démonstration a servi à quelque chose : elle a révélé \textbf{deux
défauts réels}, corrigés depuis. La réconciliation détruisait les machines des
autres nœuds, faute de savoir à qui elles appartenaient --- les machines
portent maintenant un label \cmd{insacloud.worker}. Et la reprise échouait sur
une collision de nom, parce que Docker refuse deux conteneurs homonymes :
\cmd{liberer\_le\_nom()} fait le ménage sur le nœud de repli avant de créer.
Ce second cas se produit aussi sur un vrai parc : une création qui réussit,
un enregistrement qui échoue, et la tentative suivante bute sur le reliquat.
\end{piege}

\fonc{Fonctionnalité 13 (haute disponibilité) · Test 6 de l'énoncé}

\minute{11:30 — Les six scénarios d'intégration (40 s)}

\begin{afaire}{taper}
\begin{lstlisting}[language=bash]
cd 2_webapp
ARRET_NOEUD='docker stop insacloud-%s' \
DEMARRAGE_NOEUD='docker start insacloud-%s' \
  ./tests/integration.sh http://127.0.0.1:8088
\end{lstlisting}
\end{afaire}

\begin{adire}
L'énoncé liste six tests. Ils sont scriptés et rejouables : infrastructure,
configuration, application, location avec connexion SSH réelle, expiration, et
panne d'un nœud. \textbf{Les six passent} --- vingt et une vérifications
réussies, aucun échec, une seule ignorée : HSTS, parce que la pile locale est
en HTTP et que l'en-tête n'aurait aucun sens sans TLS.

\medskip
Le quatrième ouvre une \textbf{vraie session SSH} dans la machine louée, et le
sixième est celui qu'on vient de faire à la main.

\medskip
À côté, \textbf{168 tests unitaires} en vingt secondes. Et ils ne lancent aucun
conteneur : \cmd{run\_docker} est remplacé par une fonction espion. On vérifie
donc que la commande imposée part bien telle quelle, sur un agent d'intégration
qui n'a même pas Docker.
\end{adire}

\lstinputlisting[basicstyle=\ttfamily\scriptsize]{sorties/tests.txt}

\medskip
Et le compte rendu des six scénarios, tel qu'il est sorti sur la plateforme en
fonctionnement :

\lstinputlisting[basicstyle=\ttfamily\tiny]{sorties/integration.txt}

\fonc{Fonctionnalité 16 (tests d'intégration)}

\minute{12:10 — La sécurité DevSecOps (40 s)}

\begin{afaire}{montrer}
L'onglet GitHub Actions, la liste des travaux. Puis, si on a le temps, le
durcissement lu sur une machine réelle :
\end{afaire}

\lstinputlisting[basicstyle=\ttfamily\scriptsize]{sorties/durcissement.txt}

\begin{adire}
Le « Sec » de DevSecOps, c'est que la sécurité est un contrôle
\textbf{automatique}, pas un audit de fin de projet. Six familles :

\medskip
\textbf{bandit} cherche les motifs dangereux dans le code : injection, secret en
dur, appel système construit depuis une entrée utilisateur.
\textbf{pip-audit} confronte nos dépendances aux bases de CVE.
\textbf{gitleaks} cherche les secrets commis par erreur, \textbf{y compris dans
l'historique} --- effacer un secret dans un commit suivant ne l'efface pas.
\textbf{nuclei} et un script maison attaquent l'application \textbf{en marche}.
\textbf{Trivy} scanne les images : aucune faille corrigeable HIGH ou CRITICAL ne
passe. Et \textbf{Syft} produit le SBOM, l'inventaire de ce qu'on livre.
\end{adire}

\fonc{Fonctionnalité 14 (sécurité DevSecOps)}

\minute{12:50 — La chaîne, et ce qu'elle a trouvé (40 s)}

\lstinputlisting[basicstyle=\ttfamily\scriptsize]{sorties/ci.txt}

\begin{adire}
Onze travaux à chaque envoi. Et ils \textbf{bloquent} : aucune image n'atteint
le registre si l'un échoue. Ce n'est pas décoratif --- nos propres images ont
été refusées par notre propre scan, et il a fallu changer de base.

\medskip
Dix défauts ont été trouvés en \textbf{faisant tourner} le projet, pas en le
relisant : le Faucheur qui se supprimait lui-même, le mot de passe journalisé en
clair, une course entre les quatre processus Gunicorn, une dépendance qui
n'existait pas pour le Python de la cible. Chacun a maintenant son test ou son
garde-fou --- la chaîne vérifie même sa propre syntaxe shell avant de démarrer.
\end{adire}

\fonc{Fonctionnalité 15 (CI/CD)}

\begin{center}\small\itshape
Transition : « \textit{[Prénom 1]} conclut. »
\end{center}

% =============================================================================
\chapter{Personne 1 — Conformité, ajouts, conclusion}
\markboth{PERSONNE 1 — 13:30 À 15:00}{}

\minute{13:30 — La grille (50 s)}

\begin{adire}
Les dix-sept fonctionnalités de l'énoncé sont traitées. Rapidement, l'outil
employé pour chacune : Git et un environnement épinglé pour la préparation ;
Flask et Gunicorn pour l'API ; SQLAlchemy et Alembic pour la base --- sept
tables, migrations versionnées ; Docker et supervisord pour la
conteneurisation ; Vagrant pour l'infrastructure ; Ansible pour
l'automatisation, neuf rôles idempotents ; un agent maison pour
l'enregistrement ; un ordonnanceur qui prend le nœud joignable le moins chargé ;
le Faucheur pour l'expiration et la reprise ; bandit, pip-audit, gitleaks,
nuclei, Trivy et Syft pour la sécurité ; GitHub Actions pour la chaîne ; pytest
pour les tests.

\medskip
Et pour la documentation : le README, la doc de l'API, celle de la base, quatre
schémas --- et un \textbf{guide technique de 93 pages} qui explique tout le
projet depuis les notions de base.
\end{adire}

\begin{center}
\scriptsize
\begin{tabularx}{\textwidth}{@{}clX@{}}
\toprule
\textbf{\#} & \textbf{Fonctionnalité} & \textbf{Outil employé} \\
\midrule
1 & Préparation, environnement & Git, environnement virtuel, dépendances épinglées \\
2 & API Web Flask & Flask, Jinja2, Gunicorn (4 ouvriers), nginx \\
3 & Base de données & SQLAlchemy 2.0, Alembic, PostgreSQL ou SQLite \\
4 & Conteneurisation Docker & 4 Dockerfiles, 6 images, supervisord, durcissement \\
5 & Infrastructure Vagrant & 4 VM, réseau privé, 4 providers \\
6 & Automatisation Ansible & 9 rôles, 4 playbooks, Vault AES-256, idempotence \\
7 & Worker Agent & agent stdlib, inscription + battement de cœur \\
8 & Ordonnanceur & nœud joignable le moins chargé, départage déterministe \\
9 & Service de location & 1 à 120 min, prolongation, arrêt \\
10 & Déploiement des instances & \cmd{subprocess} + la commande imposée \\
11 & Accès SSH & mot de passe ou clé, + ttyd + noVNC \\
12 & Expiration & le Faucheur, toutes les 10 s, avec réconciliation \\
13 & Haute disponibilité & battement de cœur, \cmd{OFFLINE}, recréation ailleurs \\
14 & Sécurité DevSecOps & bandit, pip-audit, gitleaks, nuclei, Trivy, Syft, Fernet \\
15 & CI/CD & GitHub Actions, 11 travaux, ghcr.io, SBOM \\
16 & Tests & 147 unitaires + 6 scénarios de bout en bout \\
17 & Documentation & README, API, base, 4 schémas, guide de 93 pages \\
\bottomrule
\end{tabularx}
\end{center}

\fonc{Fonctionnalité 1 et 17}

\minute{14:20 — Ce qu'on a ajouté (30 s)}

\begin{adire}
Sept choses au-delà de ce qui était demandé : un parc \textbf{distribué à trois
workers} alors qu'un seul hôte suffisait ; le \textbf{terminal web} et le
\textbf{bureau graphique} en plus du SSH ; un \textbf{coffre à secrets} qui
réaffiche le mot de passe après re-authentification ; un \textbf{watchdog}
systemd qui répare le service tout seul ; un \textbf{inventaire dynamique} qui
interroge Vagrant, ce qui rend le déploiement indépendant du provider ; la
\textbf{sauvegarde du travail} avant destruction ; et des \textbf{schémas
engendrés depuis leur description}, que la chaîne régénère et refuse de laisser
vieillir.
\end{adire}

\minute{14:50 — Fermer (10 s)}

\begin{adire}
Le projet tourne, il est testé, et il se redéploie en une commande. Nous sommes
à votre disposition pour les questions.
\end{adire}

% =============================================================================
\chapter{Les questions probables}
\markboth{QUESTIONS}{}

\begin{center}
\small
\begin{tabularx}{\textwidth}{@{}clX@{}}
\toprule
\textbf{Qui} & \textbf{Question} & \textbf{Réponse} \\
\midrule
P1 & Pourquoi des conteneurs et pas des VM ? &
Démarrage sous la seconde, pas de noyau invité à charger. Et l'énoncé le
tranche : « des conteneurs dans notre cas ». \\
\addlinespace
P1 & Pourquoi SSH plutôt que l'API Docker ? &
Qui parle au démon Docker d'une machine en est root. Avec
\cmd{docker -H ssh://}, aucun port Docker n'est ouvert. \\
\addlinespace
P2 & Où est stocké le mot de passe root ? &
Chiffré avec Fernet, clé distincte de celle des sessions, rangée dans le Vault.
Jamais dans un cookie --- une session Flask est signée, pas chiffrée. \\
\addlinespace
P2 & Et si deux utilisateurs demandent le même port ? &
Index unique partiel en base, limité aux machines actives : le port redevient
libre dès que la machine est détruite. \\
\addlinespace
P2 & Pourquoi supervisord dans le conteneur ? &
Un conteneur exécute un processus ; il en faut cinq en mode bureau. Et le PID 1
doit adopter les orphelins, ce qu'un shell ne fait pas. \\
\addlinespace
P3 & La reprise rend-elle les fichiers ? &
Non, et c'est assumé. Ils vivaient dans la couche inscriptible du conteneur, sur
un nœud injoignable. Sans stockage partagé, seul le service repart. \\
\addlinespace
P3 & Comment ajoute-t-on un worker ? &
On l'allume. Aucune adresse n'est écrite côté contrôleur : l'agent s'inscrit
tout seul. \\
\addlinespace
P3 & Pourquoi 168 tests sans Docker ? &
\cmd{run\_docker} est remplacé par un espion : on vérifie que la commande
imposée part bien, sur un agent qui n'a pas Docker. \\
\addlinespace
P1 & Combien de temps pour tout remonter ? &
\cmd{vagrant up} puis un playbook. Tout est décrit en fichiers versionnés. \\
\bottomrule
\end{tabularx}
\end{center}

% =============================================================================
\chapter{Si quelque chose casse en direct}
\markboth{EN CAS DE PANNE}{}

\begin{retenir}
Une panne commentée vaut mieux qu'un silence. Si quelque chose ne marche pas,
on dit à voix haute ce qu'on regarde et pourquoi --- c'est aussi une
démonstration de maîtrise.
\end{retenir}

\begin{center}
\small
\begin{tabularx}{\textwidth}{@{}XX@{}}
\toprule
\textbf{Symptôme} & \textbf{Quoi faire, à voix haute} \\
\midrule
Le site ne répond pas &
\cmd{docker compose -p insacloud ps}, puis \cmd{logs web --tail 20}. Commenter
ce qu'on lit. \\
\addlinespace
Le mode bureau n'apparaît pas &
Les images \cmd{*\_desktop} ne sont pas construites --- et c'est volontaire : le
catalogue ne propose pas ce qu'il ne peut pas livrer. \\
\addlinespace
\cmd{port is already allocated} &
Un conteneur d'un essai précédent traîne :
\cmd{docker ps -a --filter "label=insacloud.role=rented-machine"}. \\
\addlinespace
La machine met du temps à répondre en SSH &
En parler : Docker publie le port avant que \cmd{sshd} écoute. C'est pour ça que
nos tests réessaient au lieu d'attendre une durée fixe. \\
\addlinespace
Le Faucheur n'a pas détruit la machine &
Il passe toutes les dix secondes. Montrer \cmd{docker logs insacloud-faucheur}
et attendre le cycle. \\
\bottomrule
\end{tabularx}
\end{center}

\section*{Ce qu'il ne faut pas faire}

\begin{itemize}
\item \textbf{Ne pas lire ce script.} Le relire trois fois, puis parler.
\item \textbf{Ne pas enchaîner les commandes en silence.} Pendant que ça charge,
  on explique.
\item \textbf{Ne pas promettre ce qu'on ne montrera pas.} Dire d'emblée
  qu'aucun hyperviseur ne tourne sur le poste, et que les deux nœuds sont des
  conteneurs. La bascule, elle, se fait en direct.
\item \textbf{Ne pas dépasser.} En cas de retard, Personne 3 coupe les six
  scénarios d'intégration (11:30) et garde le Test 6 --- c'est lui qui compte.
\end{itemize}

\vfill
\begin{center}
\color{gristexte}\small
Dépôt : \url{https://github.com/aziz3r/insacloud}\\
Les captures de ce document viennent de la plateforme telle qu'elle tourne, et
les sorties reproduites sont les vraies.
\end{center}

\end{document}
